%% ============================================================
%% definitions/finality-definitions.tex
%% ============================================================
%%
%% Canonical finality state-machine and certificate-composition
%% definitions for Quasar Horizon. Single source of truth.
%% ============================================================

\subsection*{Quasar finality states}

\begin{definition}[Finality states]
\label{def:finality-states}
A bundle root $\mathsf{NebulaRoot}_b$ progresses through the finality
state lattice:
\begin{align*}
\mathsf{Observed} &\;\subset\; \mathsf{Beam\text{-}final}
   \;\subset\; \mathsf{PQ\text{-}attested} \\
&\;\subset\; \mathsf{Pulse\text{-}sealed}
   \;\subset\; \mathsf{Horizon\text{-}final}
\end{align*}
where each state requires the previous plus the additional certificate
specified in Definition~\ref{def:horizon-cert}.
\end{definition}

\begin{definition}[Horizon certificate]
\label{def:horizon-cert}
$\mathsf{HorizonCertificate} := (\mathsf{Beam}, \mathsf{MLDSA},
\mathsf{Pulse}, \mathsf{NebulaRoot}, \mathsf{KeyEraID},
\mathsf{GroupID}, \mathsf{Generation}, \mathsf{HashSuiteID})$
where:
\begin{itemize}[nosep]
  \item $\mathsf{Beam} = \sigma_{\textsc{bls}}$ is a BLS12-381 aggregate
        signature over $\mathsf{QuasarTranscript}$
        (Definition~\ref{def:horizon-transcript}).
  \item $\mathsf{MLDSA}$ is a cert set: signer bitmap plus the per-validator
        FIPS-204 ML-DSA-65 signatures, OR a Z-Chain Groth16 rollup
        witness compressing the verification of the cert set.
  \item $\mathsf{Pulse} = \sigma_{\textsc{pulsar}}$ is a Pulsar threshold
        signature over $\mathsf{QuasarTranscript}$.
\end{itemize}
\end{definition}

\subsection*{Prism: certificate-composition verifier}

\begin{definition}[Prism]
\label{def:prism}
$\mathsf{Prism}(\mathsf{cert}) \to \{0, 1\}$ accepts a Horizon
certificate iff all six verification predicates hold:
\begin{enumerate}[nosep]
  \item Domain separation and transcript binding: every lane signs the
        same $\mathsf{QuasarTranscript}$ (Definition
        \ref{def:horizon-transcript}); $\mathsf{HashSuiteID}$,
        $\mathsf{KeyEraID}$, $\mathsf{Generation}$, $\mathsf{NebulaRoot}$
        match across lanes.
  \item BLS Beam aggregate verifies under the validator set's BLS
        public-key registry.
  \item ML-DSA cert set verifies against the signer bitmap (or Groth16
        rollup against public-inputs hash).
  \item Pulsar Pulse verifies under the active KeyEra's GroupKey via
        $\mathsf{Pulsar.Verify}$ (Definition~\ref{def:pulsar-verify}).
  \item Signer-set and validator-set hashes match the chain's epoch state.
  \item Nebula root, epoch, round linkage matches.
\end{enumerate}
$\mathsf{Horizon\text{-}final}$ is defined as
$\mathsf{Prism}(\mathsf{cert}) = 1$. It is \textbf{not}
``BLS OR Pulsar''; it is the bound conjunction of all required lanes.
\end{definition}

\subsection*{Network tier model: L1 / L2 / L3}

\begin{definition}[Lux network tiers]
\label{def:network-tiers}
A Lux-architecture chain occupies exactly one of three tiers:
\begin{description}[nosep,leftmargin=2em]
  \item[L1 (sovereign):] no $\mathsf{primary\_network\_id}$ declared.
        Own validator set, own KeyEraID lineage, own Beam, ML-DSA cert
        set, and Pulse. May or may not be declared as another chain's
        primary; this is a graph property, not a protocol distinction.
  \item[L2 (co-validating):] declares
        $\mathsf{primary\_network\_id} = N_P$.
        Activation cert binds $\mathsf{primary\_nebula\_root}$
        from $N_P$; validator set is required to be a subset of $N_P$'s
        active set at the era's Bootstrap. Not a rollup — produces its
        own Beam, ML-DSA, Pulse.
  \item[L3 (rollup / validity / app-execution):] settles or anchors to
        an L1 or L2 via succinct proofs (e.g.\ Groth16, STARK), FHE
        circuits, or fraud / validity proofs. The proof system is a
        \emph{compatibility / compression / privacy adapter}, not the PQ
        root of trust. PQ finality of an L3 inherits from the
        underlying L1/L2's Pulse + ML-DSA, never from a pairing-based
        SNARK alone.
\end{description}
A chain may transition between tiers via a Reanchor-class governance
event.
\end{definition}

\subsection*{Proof-lane classification}

\begin{definition}[Proof lane]
\label{def:proof-lane}
Distinct from the certificate lanes (Beam, ML-DSA, Pulse), Lux supports
optional \emph{succinct proof wrappers} that compress or hide
verification work:
\begin{description}[nosep,leftmargin=2em]
  \item[Groth16 wrapper (classical):] pairing-based SNARK over BLS12-381.
        Useful for compressing ML-DSA cert-set verification into a
        $\sim 192$-byte witness (Z-Chain rollup, LP-307). \textbf{Not
        post-quantum} — pairing-based assumptions break under Shor's
        algorithm. Treat as a classical compatibility/privacy layer
        only.
  \item[Future PQ-friendly wrapper:] STARK or lattice-based proof
        system replacing Groth16 for true PQ succinctness. Open
        production-research target; not a current shipping component.
\end{description}
PQ finality is provided by ML-DSA + Pulsar. The proof lane is a
wrapper, not a root of trust.
\end{definition}

\begin{remark}[Groth16 of ML-DSA]
\label{rem:groth16-not-pq}
A Groth16 proof asserting that an ML-DSA signature set verified is a
\textbf{classical succinct proof of post-quantum signature verification}.
The signature itself is PQ; the proof system wrapping the verification
is \textbf{not post-quantum}: pairing-based SNARKs over BLS12-381 are
broken under Shor's algorithm, identical to the BLS Beam. Lux
specifications must phrase this honestly: the Groth16 lane is a
compatibility / compression / privacy adapter; PQ finality lives in
ML-DSA and Pulsar.
\end{remark}
